\section{Introduction}
Large language models (LLMs) are increasingly used to simulate users, citizens,
customers, patients, students, workers, voters, and communities. In these
systems, personas are no longer only design artifacts or role-playing cues; they
are part of the infrastructure by which LLM agents are made heterogeneous. A
persona may specify demographic attributes, values, goals, history, style,
constraints, or memory, and it may be realized as a prompt, structured record,
retrieval state, or agent profile. Recent work on LLM social simulation and
role-playing agents shows how quickly this infrastructure is expanding, from
survey-derived and social-media-grounded virtual populations to personalized
role-playing agents and digital twins \citep{arxiv241203563,arxiv240418231,arxiv251121722,arxiv250714922}.
This survey therefore does not treat personas as a new idea. Instead, it focuses
on a narrower and increasingly consequential problem: how synthetic personas for
LLM agents are generated, enriched, sampled, grounded, validated, and audited.

The central risk is that persona generation can look methodologically complete
when it is only narratively complete. Many workflows move from a plausible
persona prompt directly to simulated behavior, skipping the layer where persona
attributes are specified, sampled, fused, checked, and tied to the intended
population or use. A detailed backstory can be vivid but distributionally
arbitrary; a persona set can match one-dimensional demographic marginals while
violating feasible joint structure; and a persona-conditioned LLM can match an
aggregate outcome for the wrong reason, such as stereotype reliance or prompt
artifacts. Recent critiques of ad hoc LLM-generated personas, audits of
marginal-versus-joint alignment, and studies of persona collapse make this
failure mode explicit \citep{arxiv250316527,arxiv260612433,arxiv260424698}.
Put tersely: narrative vividness is not population validity; persona validity is
not behavioral fidelity.

Existing surveys provide important adjacent maps. Role-playing-language-agent
surveys organize personas by the identity being enacted, such as demographic,
character, or individualized personas, and social-simulation surveys organize
LLM-agent systems by simulation level, architecture, application, and evaluation
\citep{arxiv240418231,arxiv241203563}. These perspectives are necessary, but
they leave a methodological question under-specified. Whereas role-playing-agent
surveys classify personas by the identity being enacted, this survey classifies
persona-generation methods by the source and construction logic of the persona
distribution and by the objective the persona set is meant to support. The gap
is not that prior work ignores personas. The gap is that persona generation
itself has not yet been synthesized as a methodological object, with its own
choices about specification, sampling, grounding, enrichment, validation,
auditing, and uncertainty.

We separate the generation problem into a pipeline of claims. A target
population, person, or scenario is first translated into a construction process;
that process yields a structured persona set; the set is then realized as a
prompt, memory, narrative, or agent state; the LLM enacts that conditioning; and
the resulting behavior, response, or interaction trace is evaluated. Different
validity questions attach to different stages. \emph{Marginal validity} asks
whether one-dimensional attribute distributions match the intended target.
\emph{Joint validity} asks whether dependencies, constraints, and feasible
attribute combinations are preserved. \emph{Enactment validity} asks whether the
LLM behaves consistently with the persona after conditioning. \emph{Use-relative
fidelity} asks how much of any of these properties is required for a particular
downstream use, such as ideation, robustness testing, survey simulation, policy
analysis, training, or model evaluation. A persona set can be coherent before
enactment and still be ignored by the LLM; conversely, an LLM can reproduce some
observed response distribution from a weak persona prompt without the persona
set being a valid representation of the target population.

This survey makes four contributions. First, it introduces a two-axis taxonomy
of synthetic persona generation methods for LLM agents. The construction axis
distinguishes Authored Archetypes, Model-Generated Personas, Population-Sampled
Personas, and Trace-Grounded Personas. The objective axis distinguishes
Population Representation, Coverage and Stress Testing, Individual Fidelity,
Behavioral Calibration, Design Communication, Bias and Harm Auditing, and Agent
/ Model Evaluation. Second, it provides a conceptual framework that separates
construction source, objective, persona validity, and enactment validity. Third,
it synthesizes evidence and failure modes from a curated corpus of recent work,
including the construction-axis and objective-axis coding files used to develop
the taxonomy. Fourth, it derives reporting and evaluation guidance for future
persona-generation studies, with particular attention to joint structure,
grounding claims, representational harms, and the distinction between persona
quality and LLM behavior.

The rest of the paper proceeds as follows. We first define the survey scope and
search process, then position persona generation against related work on
role-playing agents, social simulation, synthetic survey data, and HCI personas.
We then present the taxonomy, discuss persona-space design and sampling
identification, review evaluation practices, and close with ethical
considerations and open problems.
